% !TEX root = main.tex
% Kết luận (5% tổng điểm, 1-2 trang, không đánh số chương)

\phantomsection
\addcontentsline{toc}{chapter}{KẾT LUẬN}
\chapter*{KẾT LUẬN}

\begin{spacing}{1.2}

% KẾT LUẬN không đánh số chương: các mục đánh số 1., 2., 3., 4.
\setcounter{section}{0}
\renewcommand{\thesection}{\arabic{section}}
\renewcommand{\theHsection}{KL.\arabic{section}}
\setcounter{table}{0}
\renewcommand{\thetable}{KL\arabic{table}}
\renewcommand{\theHtable}{KL.\arabic{table}}

\begin{huongdan}
Trọng số: 5\% tổng điểm. Độ dài khuyến nghị: 1--2 trang. Phần này không phải là tóm tắt lại toàn bộ nội dung đã viết, vì người đọc đã đọc nội dung đó rồi. Đây là nơi nhìn lại toàn bộ hành trình nghiên cứu với góc nhìn tổng thể, tự đánh giá trung thực, và chỉ ra con đường tiếp theo.
\end{huongdan}

\section{Tóm tắt kết quả đạt được}
\label{sec:kl_tom_tat}

\begin{huongdan}
Không quá nửa trang. Tóm tắt mỗi chương bằng hai đến ba câu: Chương 1 đã hệ thống hóa được gì; Chương 2 xác định vấn đề gì và tìm thấy Gap nào; Chương 3 đề xuất giải pháp gì và dự kiến kết quả ra sao. Dùng ngôn ngữ kết nối giữa các chương, không liệt kê rời rạc.
\end{huongdan}

Nội dung của phần này.

\section{Tự đánh giá tính khả thi}
\label{sec:kl_kha_thi}

\begin{huongdan}
Đánh giá trung thực bốn yếu tố: dữ liệu, kỹ năng, ngân sách, thời gian. Không nói "tất cả đều ổn" nếu có rủi ro thực sự.
\end{huongdan}

Nội dung của phần này. Kết quả tự đánh giá được tổng hợp tại Bảng~\ref{tab:kha_thi}.

\begin{table}[htpb]
\centering
\caption{Tự đánh giá tính khả thi của đề cương}
\label{tab:kha_thi}
\renewcommand{\arraystretch}{1.3}
{\fontsize{12}{14}\selectfont
\begin{tabularx}{\textwidth}{|>{\raggedright\arraybackslash}p{3.5cm}|X|>{\centering\arraybackslash}p{3.5cm}|}
\hline
\textbf{Yếu tố} & \textbf{Tự đánh giá (Học viên điền)} & \textbf{Mức độ sẵn sàng} \\
\hline
Dữ liệu & [Đã được cam kết cấp phép / Chưa xác nhận / Cần phương án dự phòng] & [Cao / TB / Thấp] \\
\hline
Kỹ năng kỹ thuật & [Đã thành thạo / Cần học thêm X / Cần hỗ trợ từ Y] & [Cao / TB / Thấp] \\
\hline
Ngân sách & [Đã có nguồn / Cần xin thêm / Dùng free tier đủ] & [Cao / TB / Thấp] \\
\hline
Thời gian & [Khả thi / Có rủi ro trễ GĐ X / Cần điều chỉnh GĐ Y] & [Cao / TB / Thấp] \\
\hline
\end{tabularx}
}
\end{table}

\section{Những hạn chế của đề cương hiện tại}
\label{sec:kl_han_che}

\begin{huongdan}
Nêu thẳng hai đến ba hạn chế của đề cương (phạm vi giới hạn, giả định chưa kiểm chứng...); giải thích tại sao các hạn chế này chấp nhận được trong bối cảnh tiểu luận thạc sĩ. Đây không phải điểm trừ, vì người chấm đánh giá cao học viên nhận thức được giới hạn của mình.
\end{huongdan}

Nội dung của phần này.

\section{Bài học kinh nghiệm và hướng nghiên cứu tương lai}
\label{sec:kl_bai_hoc}

\begin{huongdan}
Nêu bài học về phương pháp luận; bài học về công cụ (Mendeley/Zotero); và hai đến ba hướng nghiên cứu tương lai cụ thể, liên kết với mục~\ref{subsec:dong_gop_mo_rong}.
\end{huongdan}

Nội dung của phần này.

\end{spacing}
